\documentclass[10pt,a4paper]{amsart}
\usepackage{amsmath,amssymb,amsthm,graphicx}
\usepackage{verbatim}
\newenvironment{reference}{\comment}{\endcomment}
%\newenvironment{reference}{}{}
\newenvironment{slogan}{\comment}{\endcomment}
\newenvironment{history}{\comment}{\endcomment}
% For commutative diagrams we use Xy-pic
\usepackage[all]{xy}

% We use 2cell for 2-commutative diagrams.
\xyoption{2cell}
% Theorem environments.
%
\theoremstyle{plain}
\newtheorem{theorem}[subsection]{Theorem}
\newtheorem{proposition}[subsection]{Proposition}
\newtheorem{lemma}[subsection]{Lemma}

\theoremstyle{definition}
\newtheorem{definition}[subsection]{Definition}
\newtheorem{example}[subsection]{Example}
\newtheorem{exercise}[subsection]{Exercise}
\newtheorem{situation}[subsection]{Situation}

\theoremstyle{remark}
\newtheorem{remark}[subsection]{Remark}
\newtheorem{remarks}[subsection]{Remarks}

\numberwithin{equation}{subsection}

% Macros
%
\def\lim{\mathop{\mathrm{lim}}\nolimits}
\def\colim{\mathop{\mathrm{colim}}\nolimits}
\def\Spec{\mathop{\mathrm{Spec}}}
\def\Hom{\mathop{\mathrm{Hom}}\nolimits}
\def\Ext{\mathop{\mathrm{Ext}}\nolimits}
\def\SheafHom{\mathop{\mathcal{H}\!\mathit{om}}\nolimits}
\def\SheafExt{\mathop{\mathcal{E}\!\mathit{xt}}\nolimits}
\def\Sch{\mathit{Sch}}
\def\Mor{\mathop{\mathrm{Mor}}\nolimits}
\def\Ob{\mathop{\mathrm{Ob}}\nolimits}
\def\Sh{\mathop{\mathit{Sh}}\nolimits}
\def\NL{\mathop{N\!L}\nolimits}
\def\CH{\mathop{\mathrm{CH}}\nolimits}
\def\proetale{{pro\text{-}\acute{e}tale}}
\def\etale{{\acute{e}tale}}
\def\QCoh{\mathit{QCoh}}
\def\Ker{\mathop{\mathrm{Ker}}}
\def\Im{\mathop{\mathrm{Im}}}
\def\Coker{\mathop{\mathrm{Coker}}}
\def\Coim{\mathop{\mathrm{Coim}}}

% Boxtimes
%
\DeclareMathSymbol{\boxtimes}{\mathbin}{AMSa}{"02}

%
% Macros for moduli stacks/spaces
%
\def\QCohstack{\mathcal{QC}\!\mathit{oh}}
\def\Cohstack{\mathcal{C}\!\mathit{oh}}
\def\Spacesstack{\mathcal{S}\!\mathit{paces}}
\def\Quotfunctor{\mathrm{Quot}}
\def\Hilbfunctor{\mathrm{Hilb}}
\def\Curvesstack{\mathcal{C}\!\mathit{urves}}
\def\Polarizedstack{\mathcal{P}\!\mathit{olarized}}
\def\Complexesstack{\mathcal{C}\!\mathit{omplexes}}
% \Pic is the operator that assigns to X its picard group, usage \Pic(X)
% \Picardstack_{X/B} denotes the Picard stack of X over B
% \Picardfunctor_{X/B} denotes the Picard functor of X over B
\def\Pic{\mathop{\mathrm{Pic}}\nolimits}
\def\Picardstack{\mathcal{P}\!\mathit{ic}}
\def\Picardfunctor{\mathrm{Pic}}
\def\Deformationcategory{\mathcal{D}\!\mathit{ef}}
\usepackage{cleveref}
\crefname{originallemma}{original lemma}{original lemmas}
\Crefname{originallemma}{Original lemma}{Original lemmas}
\makeatletter
\newlabel{lemma-compute-for-effective-Cartier-algebraic}{{lemma-compute-for-effective-Cartier-algebraic}{source}{}{sourcepointer-1}{}}
\newlabel{lemma-compute-for-effective-Cartier-algebraic@cref}{{[originallemma][0][]lemma-compute-for-effective-Cartier-algebraic}{[1][0][]source}{}{}{}}
\makeatother
\crefname{originallemma}{original lemma}{original lemmas}
\Crefname{originallemma}{Original lemma}{Original lemmas}
\makeatletter
\newlabel{lemma-divide-by-nonzerodivisor}{{lemma-divide-by-nonzerodivisor}{source}{}{sourcepointer-2}{}}
\newlabel{lemma-divide-by-nonzerodivisor@cref}{{[originallemma][0][]lemma-divide-by-nonzerodivisor}{[1][0][]source}{}{}{}}
\makeatother
\crefname{originallemma}{original lemma}{original lemmas}
\Crefname{originallemma}{Original lemma}{Original lemmas}
\makeatletter
\newlabel{lemma-gorenstein-ext}{{lemma-gorenstein-ext}{source}{}{sourcepointer-3}{}}
\newlabel{lemma-gorenstein-ext@cref}{{[originallemma][0][]lemma-gorenstein-ext}{[1][0][]source}{}{}{}}
\makeatother
\crefname{originallemma}{original lemma}{original lemmas}
\Crefname{originallemma}{Original lemma}{Original lemmas}
\makeatletter
\newlabel{lemma-right-adjoint}{{lemma-right-adjoint}{source}{}{sourcepointer-4}{}}
\newlabel{lemma-right-adjoint@cref}{{[originallemma][0][]lemma-right-adjoint}{[1][0][]source}{}{}{}}
\makeatother
\title{Being Gorenstein is preserved and reflected by dividing by a nonzerodivisor}
\author{The Stacks Project Authors}
\date{stacks-project:dualizing.tex:local-snapshot-20261012; GNU FDL 1.2 or later}
\begin{document}
\maketitle
\noindent\textit{Editorial note.} All delivered mathematical text is the author's own, copied from the pinned Stacks Project chapter source sources/stacks-dualizing/pinned-original-source.tex; no statement or proof step was rewritten, reordered or removed.\par\medskip
\noindent\textit{Editorial note.} Preamble provenance: the Stacks Project keeps one single shared preamble for every chapter; this card takes it directly from the staged sources/stacks-dualizing/preamble.tex as three verbatim spans (lines 8-12, 14-18 and 154-217), so no derived file was written and nothing was added or redefined. The wrapper class and the base packages are supplied by the assembler.\par\medskip
\noindent\textit{Editorial note.} Bibliography provenance: neither the counted statement nor its proof contains a citation (this lane rescans the whole span before writing the card), so the card carries no bibliography and no entry was reconstructed.\par\medskip
\noindent\textit{Editorial note.} Selection provenance: this card was opened under the file-based control contract (CONTROL.md), with the mechanical pre-card gate run on its statement and proof interval before the card was written.\par\medskip
\noindent\textit{Allowed standard prerequisites (editorial).}\begin{itemize}
\item The preservation of the Gorenstein property under dividing by a nonzerodivisor. Invoked at the start of the counted proof (source lines 4117-4118). Gives the forward implication.
\item The vanishing criterion for the Gorenstein property through the Ext modules of the residue field. Invoked inside the counted proof (source lines 4119-4121 and 4131-4132). Translates between Gorensteinness and the eventual vanishing of those Ext modules in both directions.
\item The adjunction between restriction and the internal hom along the quotient map. Invoked inside the counted proof (source lines 4122-4123). Computes the Ext modules of the residue field over the source in terms of those over the quotient.
\item The computation of the internal hom out of the quotient as the quotient shifted by one. Invoked inside the counted proof (source lines 4129-4130). Turns the computed internal hom into the shift of the quotient, matching the criterion.
\end{itemize}\medskip
\ifdefined\hypertarget\hypertarget{sourcepointer-1}{}\fi
\noindent\textit{Source reference.} Original lemma lemma-compute-for-effective-Cartier-algebraic: sources/stacks-dualizing/pinned-original-source.tex, source lines 2621--2621. Navigation pointer only: the counted text refers to the Stacks item lemma-compute-for-effective-Cartier-algebraic, which is an existing labeled result of this same chapter and is not part of this card. The pointer carries the target's own label as its number and its definition line as its span; its content is not reproduced here.\par\medskip
\ifdefined\hypertarget\hypertarget{sourcepointer-2}{}\fi
\noindent\textit{Source reference.} Original lemma lemma-divide-by-nonzerodivisor: sources/stacks-dualizing/pinned-original-source.tex, source lines 3383--3383. Navigation pointer only: the counted text refers to the Stacks item lemma-divide-by-nonzerodivisor, which is an existing labeled result of this same chapter and is not part of this card. The pointer carries the target's own label as its number and its definition line as its span; its content is not reproduced here.\par\medskip
\ifdefined\hypertarget\hypertarget{sourcepointer-3}{}\fi
\noindent\textit{Source reference.} Original lemma lemma-gorenstein-ext: sources/stacks-dualizing/pinned-original-source.tex, source lines 4063--4063. Navigation pointer only: the counted text refers to the Stacks item lemma-gorenstein-ext, which is an existing labeled result of this same chapter and is not part of this card. The pointer carries the target's own label as its number and its definition line as its span; its content is not reproduced here.\par\medskip
\ifdefined\hypertarget\hypertarget{sourcepointer-4}{}\fi
\noindent\textit{Source reference.} Original lemma lemma-right-adjoint: sources/stacks-dualizing/pinned-original-source.tex, source lines 2428--2428. Navigation pointer only: the counted text refers to the Stacks item lemma-right-adjoint, which is an existing labeled result of this same chapter and is not part of this card. The pointer carries the target's own label as its number and its definition line as its span; its content is not reproduced here.\par\medskip
\begin{lemma}
\label{lemma-gorenstein-divide-by-nonzerodivisor}
Let $(A, \mathfrak m, \kappa)$
be a Noetherian local ring. Let $f \in \mathfrak m$ be a
nonzerodivisor. Set $B = A/(f)$. Then $A$ is Gorenstein if and
only if $B$ is Gorenstein.
\end{lemma}

\begin{proof}
If $A$ is Gorenstein, then $B$ is Gorenstein by
Lemma \ref{lemma-divide-by-nonzerodivisor}.
Conversely, suppose that $B$ is Gorenstein. Then
$\Ext^i_B(\kappa, B)$ is zero for $i \gg 0$
(Lemma \ref{lemma-gorenstein-ext}).
Recall that $R\Hom(B, -) : D(A) \to D(B)$ is a right adjoint
to restriction (Lemma \ref{lemma-right-adjoint}).
Hence
$$
R\Hom_A(\kappa, A) = R\Hom_B(\kappa, R\Hom(B, A)) =
R\Hom_B(\kappa, B[1])
$$
The final equality by direct computation or by
Lemma \ref{lemma-compute-for-effective-Cartier-algebraic}.
Thus we see that $\Ext^i_A(\kappa, A)$ is zero for
$i \gg 0$ and $A$ is Gorenstein (Lemma \ref{lemma-gorenstein-ext}).
\end{proof}

\end{document}
